\section{Exact local and center formulas}
\label{app:formulas}

This appendix records the formulas used by the propagation arguments in
Section~\ref{sec:exact-demand}.  Every radical below denotes its nonnegative
value, and every algebraic root is restricted to the stated geometric
contact interval.  This last convention is important: an unrestricted root
of a squared support equation need not describe an enclosing triangle.

\subsection{The local admissible set}

Put
\[
 u=\left(\frac12,\frac{\sqrt3}{2}\right),\qquad
 v=\left(\frac12,-\frac{\sqrt3}{2}\right)
\]
and, for $0\le a,b,c\le1$, let
\[
 K(a,b,c)=\conv\{0,au,bv,c(u+v)\}.
\]
Write $L_{\min}(a,b,c)$ for the least side length of an equilateral
triangle containing this set, and let
\[
 \mathcal A=\{(a,b,c)\in[0,1]^3:L_{\min}(a,b,c)\le1\}.
\]
Set $s=a+b$, $D=(a^2+ab+b^2)^{1/2}$, and
\[
 R_a=(a^2-ac+c^2)^{1/2},\qquad
 R_b=(b^2-bc+c^2)^{1/2}.
\]
When $s>0$ and $cs>ab$, define
\[
 L_{OA}=b+\max\{a,c\},\qquad
 L_{BO}=a+\max\{b,c\},
\]
\[
 L_{AC}=
 \begin{cases}
 (a^2+bc)/R_a,&a\ge c,\\
 c(a+b)/R_a,&a<c\le a+b,\\
 c^2/R_a,&c>a+b,
 \end{cases}
\]
and let $L_{CB}$ be the same formula with $a$ and $b$ interchanged.
The support-function computation in the proof of
Proposition~\ref{prop:exact-admissible-set} gives
\begin{equation}
 L_{\min}(a,b,c)=
 \begin{cases}
 c,&a=b=0,\\
 D,&s>0\text{ and }cs\le ab,\\
 \min\{L_{OA},L_{BO},L_{AC},L_{CB}\},
     &s>0\text{ and }cs>ab.
 \end{cases}
 \label{eq:Lmin-support-formula}
\end{equation}

For the polynomial form, put
\[
 m=\min\{a,b\},\quad M=\max\{a,b\},\quad
 d=a^2+ab+b^2,
 \quad q=s^4-s^2+ab,
\]
\[
 F_L=c^4-c^2+mc-m^2,
\]
\[
 F_T=(s^2-1)c^2+Mc-M^2,
\]
\[
 F_S=(m^2-1)c^2+(2mM^2+M)c+(M^4-M^2).
\]
Then
\begin{equation}
 \mathcal A=\mathcal A_L\cup\mathcal A_T\cup\mathcal A_S,
 \label{eq:admissible-cells}
\end{equation}
where
\[
 \begin{aligned}
 \mathcal A_L={}&\{d\le1, s\le1, q\le0, F_L\le0\},\\
 \mathcal A_T={}&\{d\le1, s\le1, q\ge0, F_T\le0, c\le2M\},\\
 \mathcal A_S={}&\{d\le1, s>1, F_S\le0, c\le\tfrac12\}.
 \end{aligned}
\]
All three sets are understood to lie in $[0,1]^3$.  The selector
$c\le2M$ removes the spurious upper sheet of $F_T\le0$, and
$c\le1/2$ selects the lower component in the supercritical cell.

\subsection{The exact outgoing map}

For $a,c\in[0,1]$, define
\[
 B_c(a)=\max\{b\in[0,1]:(a,b,c)\in\mathcal A\},
 \qquad
 \beta(a)=\frac{-a+\sqrt{4-3a^2}}2.
\]
The fiber is the whole interval $[0,B_c(a)]$.  We now give the finite
piecewise-algebraic formula used in the proof.  The following catalog assumes
$0<a,c<1$ and writes $\beta=\beta(a)$; an inapplicable entry is assigned
$-\infty$.

The inner and axial candidates are
\[
 \beta_{\rm in}=
 \begin{cases}\beta,&c(a+\beta)\le a\beta,\\-\infty,&\text{otherwise},\end{cases}
 \qquad
 \beta_{OA}=\min\{\beta,1-\max\{a,c\}\},
\]
\[
 \beta_{BO}=
 \begin{cases}\min\{\beta,1-a\},&a+c\le1,\\-\infty,&a+c>1.
 \end{cases}
\]
For the two $AC$ contacts set
\[
 U_A=\begin{cases}+\infty,&a=c,\\ac/(a-c),&a>c,
 \end{cases}
\]
and
\[
 \beta_{AC}^{(1)}=
 \begin{cases}
 \min\{\beta,U_A,(R_a-a^2)/c\},
 &a\ge c\text{ and this minimum is nonnegative},\\
 -\infty,&\text{otherwise},
 \end{cases}
\]
\[
 \beta_{AC}^{(2)}=
 \begin{cases}
 \min\{\beta,R_a/c-a\},&a<c\text{ and }R_a\ge c^2,\\
 -\infty,&\text{otherwise}.
 \end{cases}
\]

There are three $CB$ contacts.  Define
\[
 P_{a,c}(b)=b^4+(2ac-1)b^2+cb+(a^2-1)c^2
\]
and
\[
 U_H=\begin{cases}
 \beta,&a\le c,\\
 \min\{\beta,ac/(a-c)\},&a>c.
 \end{cases}
\]
If $c\le U_H$, let $\rho_H$ denote the largest root of $P_{a,c}$ in
$[c,U_H]$, when such a root exists.  Then
\[
 \beta_{CB}^{(1)}=
 \begin{cases}
 U_H,&c\le U_H\text{ and }P_{a,c}(U_H)\le0,\\
 \rho_H,&c\le U_H, P_{a,c}(U_H)>0,
              \text{ and }\rho_H\text{ exists},\\
 -\infty,&\text{otherwise}.
 \end{cases}
\]
Next put
\[
 L_M=\max\{0,c-a\},\quad U_M=\min\{c,\beta\},
\]
\[
 Q_{a,c}(b)=(c^2-1)b^2+(2ac^2+c)b+(a^2-1)c^2,
\]
\[
 \Delta_M=(2ac^2+c)^2-4(1-c^2)(1-a^2)c^2,
 \quad
 r_\pm=\frac{2ac^2+c\pm\sqrt{\Delta_M}}{2(1-c^2)}
\]
when $\Delta_M\ge0$.  The middle contact is
\[
 \beta_{CB}^{(2)}=
 \begin{cases}
 -\infty,&L_M>U_M,\\
 U_M,&L_M\le U_M\text{ and }Q_{a,c}(U_M)\le0,\\
 r_-,&\Delta_M\ge0\text{ and }L_M\le r_-<U_M<r_+,\\
 -\infty,&\text{otherwise}.
 \end{cases}
\]
If $\Delta_M<0$, the second line applies whenever $L_M\le U_M$.
Finally set $U_L=\min\{\beta,c-a\}$ and, for $c\ge\sqrt3/2$,
\[
 \ell(c)=\frac c2\left(1-\sqrt{4c^2-3}\right),\qquad
 r(c)=\frac c2\left(1+\sqrt{4c^2-3}\right).
\]
Then
\[
 \beta_{CB}^{(3)}=
 \begin{cases}
 -\infty,&U_L<0,\\
 U_L,&U_L\ge0\text{ and }c<\sqrt3/2,\\
 U_L,&c\ge\sqrt3/2\text{ and }
       (U_L\le\ell(c)\text{ or }U_L\ge r(c)),\\
 \ell(c),&c\ge\sqrt3/2\text{ and }\ell(c)<U_L<r(c).
 \end{cases}
\]
With
\[
 \mathcal B(a,c)=\{\beta_{\rm in},\beta_{OA},\beta_{BO},
 \beta_{AC}^{(1)},\beta_{AC}^{(2)},
 \beta_{CB}^{(1)},\beta_{CB}^{(2)},\beta_{CB}^{(3)}\}\setminus\{-\infty\},
\]
the final formula is
\begin{equation}
 B_c(a)=
 \begin{cases}
 \beta(a),&c=0,\\
 1,&a=0, 0<c\le1,\\
 0,&a=1, 0<c\le1,\\
 0,&0<a<1, c=1,\\
 \max\mathcal B(a,c),&0<a,c<1.
 \end{cases}
 \label{eq:exact-B-map}
\end{equation}
Every entry in $\mathcal B$ is obtained from one of the exhaustive support
contacts in~\eqref{eq:Lmin-support-formula}; hence the maximum introduces no
unverified algebraic branch.

\subsection{The capped nonsupercritical map}

Define
\[
 F_c(a)=\min\{B_c(a),1-a\},\qquad G_c(a)=1-F_c(a).
\]
For $1/2<c<1$, set
\[
 h(c)=\frac{2c}{\sqrt{16c^2-3}+1},
\]
\[
 T_-(a,c)=\frac{\sqrt{a^2-ac+c^2}}c-a,
\]
and
\[
 T_+(a,c)=
 \frac{2(1-a^2)c^2}{2ac^2+c+\sqrt{\Delta(a,c)}},
 \quad
 \Delta(a,c)=(2ac^2+c)^2-4(1-c^2)(1-a^2)c^2.
\]
For $1/2<c\le\sqrt3/2$,
\begin{equation}
 F_c(a)=
 \begin{cases}
 1-a,&0\le a\le1-c,\\
 T_+(a,c),&1-c<a<h(c),\\
 h(c),&a=h(c),\\
 T_-(a,c),&h(c)<a<c,\\
 1-a,&c\le a\le1.
 \end{cases}
 \label{eq:capped-map-low}
\end{equation}
For $\sqrt3/2<c<1$,
\begin{equation}
 F_c(a)=
 \begin{cases}
 1-a,&0\le a\le1-c,\\
 T_+(a,c),&1-c<a\le\ell(c),\\
 \ell(c),&\ell(c)<a<r(c),\\
 T_-(a,c),&r(c)\le a<c,\\
 1-a,&c\le a\le1.
 \end{cases}
 \label{eq:capped-map-high}
\end{equation}
At $a=\ell(c)$ the value is $r(c)$; immediately to its right it drops to
$\ell(c)$.  Thus endpoints may not be reassigned by continuity.  The four
genuine labels are Full, $L$, $T_+$, and $T_-$.  The map is nonincreasing,
$G_c(a)\ge a$, and reflection of $\mathcal A$ gives the exact duality
\begin{equation}
 G_c(a)\le z\quad\Longleftrightarrow\quad
 G_c(1-z)\le1-a.
 \label{eq:capped-duality}
\end{equation}

\subsection{Exact center models}

We record the normalized forms after rotating and, when necessary,
reflecting so that the unique center-covered midpoint is $M_0$.

For CE1, use affine coordinates
\[
 X=V_0+b(V_1-V_0)+a(V_5-V_0)
\]
and write $T_C\cap e_{0,1}=[s,t]$.  Put
\[
 0<\lambda<1,\qquad \rho=\sqrt{1-\lambda+\lambda^2},
\]
\[
 \begin{aligned}
 F_1&=\lambda b+(1-\lambda)a-\lambda s,\\
 F_2&=-b+\lambda a+t,\\
 F_0&=(1-\lambda)b-a+\rho+\lambda s-t,
 \end{aligned}
\qquad
 T_C=\{F_0,F_1,F_2\ge0\}.
\]
At $O$, the three slacks are
\[
 C_0=\rho+\lambda s-t-\lambda,\quad
 C_1=1-\lambda s,\quad
 C_2=t+\lambda-1,\quad
 C_0+C_1+C_2=\rho.
\]
The exact closed domain is
\begin{align*}
 &0<\lambda<1,\quad0\le s<t\le1,\quad
 t+\lambda(1-s)\le\rho,\quad t\ge1-\lambda,\\
 &\lambda s\le\tfrac12,\quad t<1-\tfrac\lambda2,\quad
 t>\rho+\lambda s-\tfrac{1+\lambda}{2},\quad
 t\ge\rho-\frac{\lambda^2}{1-\lambda}s.
\end{align*}
For the closure of an open center role the first two center slacks satisfy
$C_0,C_2>0$.  The radial exits $d_i$ and complementary demands $c_i=1-d_i$
are
\begin{equation}
 \begin{aligned}
 d_0&=1-\lambda s,&c_0&=\lambda s,\\
 d_1&=C_2/\lambda,&c_1&=1-C_2/\lambda,\\
 d_2&=C_2,&c_2&=1-C_2,\\
 d_3&=\min\{C_0/\lambda,C_2/(1-\lambda)\},
 &c_3&=1-d_3,\\
 d_4&=C_0,&c_4&=1-C_0,\\
 d_5&=C_0/(1-\lambda),&c_5&=1-C_0/(1-\lambda).
 \end{aligned}
 \label{eq:CE1-exits}
\end{equation}

For CE2, parameterize the two adjacent edge traces from $V_0$ by
\[
 T_C\cap e_{5,0}=[x,u],\qquad T_C\cap e_{0,1}=[y,v]
\]
and put
\[
 S=x+y,\qquad D=\sqrt{x^2+xy+y^2}.
\]
The exact domain is
\begin{align*}
 &0<x<u<1,\qquad0<y<v<1,\\
 &(u+v)S-xy=D,\qquad uS\ge x,\qquad vS\ge y,\\
 &uS<x+\tfrac y2,\qquad vS<y+\tfrac x2.
\end{align*}
The two center inequalities become strict for the closure of an open center
role: $uS>x$ and $vS>y$.  In particular,
\begin{equation}
 u+v=\frac{D+xy}{x+y};
 \label{eq:CE2-coupling}
\end{equation}
the two edge intervals cannot be varied independently.  The exact exits are
\begin{equation}
 \begin{aligned}
 d_0&=1-xy/S,&d_1&=(vS-y)/x,&d_2&=(vS-y)/S,\\
 d_3&=\min\{(vS-y)/y,(uS-x)/x\},\\
 d_4&=(uS-x)/S,&d_5&=(uS-x)/y,
 \end{aligned}
 \label{eq:CE2-exits}
\end{equation}
and again $c_i=1-d_i$.  In both normalized center classes,
\[
 d_0\ge\frac12,\qquad d_i<\frac12\ (i\ne0),
 \qquad
 c_0\le\frac12,\qquad c_i>\frac12\ (i\ne0).
\]
